\documentclass[11pt]{article}

\newcommand{\cnum}{Math 245C}
\newcommand{\ced}{Spring 2026}
\newcommand{\ctitle}[4]{\title{\vspace{-0.5in}\cnum, \ced\\week #1: #2}\author{\vspace{-0.35in}\\#3, #4}}
\usepackage{enumitem}
\usepackage[usenames,dvipsnames,svgnames,table,hyperref]{xcolor}
\usepackage{amsmath, amsfonts}
\usepackage{graphicx} % Required for inserting images
\usepackage{float}
\usepackage{listings}
\usepackage{xcolor}
\usepackage{hyperref}
\usepackage{comment}

\renewcommand*{\theenumi}{\alph{enumi}}
\renewcommand*\labelenumi{(\theenumi)}
\renewcommand*{\theenumii}{\roman{enumii}}
\renewcommand*\labelenumii{\theenumii.}
\newcommand{\notiff}{%
  \mathrel{{\ooalign{\hidewidth$\not\phantom{"}$\hidewidth\cr$\iff$}}}}

\definecolor{codegreen}{rgb}{0,0.6,0}
\definecolor{codegray}{rgb}{0.5,0.5,0.5}
\definecolor{codepurple}{rgb}{0.58,0,0.82}
\definecolor{backcolour}{rgb}{0.95,0.95,0.92}
\definecolor{header}{HTML}{205459}
\definecolor{solution}{HTML}{204d52}

\newcommand{\solution}[1]{{{\textcolor{header}{Solution:} \textcolor{solution}{#1}}}}
\setlist[enumerate]{label=(\alph*)}

\lstdefinestyle{mystyle}{
    backgroundcolor=\color{backcolour},
    commentstyle=\color{codegreen},
    keywordstyle=\color{magenta},
    numberstyle=\tiny\color{codegray},
    stringstyle=\color{codepurple},
    basicstyle=\ttfamily\footnotesize,
    breakatwhitespace=false,
    breaklines=true,
    captionpos=b,
    keepspaces=true,
    numbers=left,
    numbersep=5pt,
    showspaces=false,
    showstringspaces=false,
    showtabs=false,
    tabsize=2
}


\begin{document}
\ctitle{\#1}{}{Asher Christian}{}
\date{}
\maketitle

\section{Monday}

\section{Lusin's Theorem}
Suppose $\mu $ is a regular Borel measure on $ \mathbb{R}^{d}$,
$A \subset \mathbb{R}^{d}$ is $\mu $-measurable and $\mu[A] < \infty$ if $f : A \to \mathbb{R}$ is $\mu$-measurable  then
for every $\epsilon > 0$ there exists $K \subset A$ compact such that $\mu[A \setminus K] < \epsilon$ 
and $f|_K$ is continuous and you can replace  $ \mathbb{R}^{d}$ by $X$ metric space complete separable.
Folland(textbook)
Distributions and weak $L^{p}$-functions.
Throughout this section $(\Omega, \Sigma, \mu)$ is a measure space Let $\nu := g_\# \mu$ be the push forward measure of $g$ where
\[
    \nu[B] = \mu[g^{-1}(B)]
\] 
you can also define
\[
\mu = g^{*}\nu
\] 
where
\[
    \mu[A] = \nu[g(A)]
\] 
example $\Lambda = (0,+\infty)$
\section{Definition}
We define $\lambda_f$ the distribution function of a $\mu$-measurable function $f : \Omega \to [-\infty,\infty]$ by $\lambda_f(\alpha) = \mu\{ \alpha < |f|\} = \mu \{\omega \in \Omega: \alpha < |f(\omega)|\}$
\[
= |f|^{-1}(\alpha, +\infty)
\] 
so if $g = |f|$ this is the  $\nu$ pullback of these intervals.
\section{Theorem}
 \begin{enumerate}
     \item $\lambda_f$ is monotone, non increasing and right continuous
     \item if $|f| \le |g|$ then  $\lambda_f \le \lambda_g$
     \item if $(|f_n|)_n$ increases to  $|f|$ as  $n \to \infty$ then $\lim_{n\to \infty}\lambda_{f_n}(\alpha) = \lambda_f(\alpha)$ for all $\alpha > 0$
     \item If $f = g + h$ then  $\lambda_f(\alpha) \le \lambda_g(\frac{\alpha}{2}) + \lambda_h(\frac{\alpha}{2})$
\end{enumerate}
\[
    \lambda_{f_1 + \dots + f_n} = \sum_{i=1}^{n} \lambda_{f_i}
\] 
if for all $i \ne j$  $\emptyset = \{|f_i| > 0\} \cap \{|f_j| > 0\}$.
Proof if  $0 < \alpha_1 < \alpha_2$ then
\[
    \{\alpha_1 < |f|\} \supset \{\alpha_2 < |f|\}
\] 
so
\[
\lambda_f(\alpha_1) \ge \lambda_f(\alpha_2)
\] 
Since $\lambda_f$ is monotone  nonincreasing to show that it is right continuous it suffices to show
that
\[
\lim_{k\to +\infty} \lambda_f(\alpha + \frac{1}{k}) = \lambda_f(\alpha)
\] 
but
\[
    \{\alpha < |f|\} = \bigcup_{k=1}^{\infty}\{\alpha + \frac{1}{k} < |f|\}
\] 
so
\[
    \mu\{\alpha < |f|\} = \lim_{k\to \infty} \mu\{\alpha + \frac{1}{k} < |f|\}
\] 
so
\[
\lambda_f(\alpha) = \lim_{k\to \infty} \lambda_f(\alpha + \frac{1}{k})
\] 
For 3 assume that $|f_n| \uparrow f$ then  $\{\alpha < |f|\}$ if for $\omega, \alpha < |f(\omega)|$ then for some $N$ all  $n > N$  $\alpha < |f_n(\omega)|$
so
\[
    \{\alpha < |f|\} = \bigcup_{n=1}^\infty \{\alpha < |f_n|\} = \bigcup_{n=1}^{\infty}A_n
\] 
and each $A_n$ are  $\mu$-measurable so we conclude that
\[
    \lambda_f(\alpha) = \lim_{n\to \infty}\lambda_{f_n}(\alpha)
\] 
assume $f = g + h$ and  $\alpha > 0$ then $\alpha < |f| \le |g| + |h|$  this implies that
\[
\frac{\alpha}{2} < |g|
\] 
or 
\[
\frac{\alpha}{2} < |h|
\] 
so
\[
    \{\alpha < |f|\} \subset \{\frac{\alpha}{2} < |g|\} \cup \{\frac{\alpha}{2} < |h|\}
\] 
so
\[
    \mu\{\alpha < |f|\} \le \mu\{\frac{\alpha}{2} < |g|\} + \mu\{\frac{\alpha}{2} < |h|\}
\] 
Remark: Since $-\lambda_f$ is monotone non decreasing and right continuous there exists a Borel measure  $\nu_f$ on  $(0,+\infty)$ such that
\[
    \nu_f(a,b] = -\lambda_f(b) - (- \lambda_f(a)) = \lambda_f(a) - \lambda_f(b) = \mu\{a < |f| \le b\} = \int_{\Omega}^{} 1_{(a,b]}(|f|) d \mu
\] 
and
\[
    \{a < |f|\} = \{a < |f| \le b\} \cup \{b < |f|\}
\] 
and so
\[
    \lambda_f(a) = \mu\{a < |f| \le b\} + \lambda_f(b)
\] 

\[
    \mu[B] = \int_{\Omega}^{}1_B d \mu
\] 
\[
    1_{(a,b]}(|f|) = 1_{|f|^{-1}(a,b]}
\] 

\section{Wednesday}

Reminder: $(\Omega, \Sigma (M), \mu)$ is a measure space and for $f : \Omega \to [-\infty, +\infty]$ $\mu$-measurable
\[
    \lambda_f(\alpha) = \mu\{\alpha < |f| \} \qquad \forall \alpha > 0
\] 
Remark (Correction) If $\lambda_f(\alpha) < +\infty$ for all $\alpha > 0$ then there exists a negative borel measure
$\nu_f$ on  $(0,+\infty)$ such that
\[
    \nu_f(a,b] = \lambda_f(b) - \lambda_f(a) \text{if $(a,b] \subset (0,+\infty)$}
\] 
we can make it a regular borel measure by setting
\[
    (1) \quad   \int_{(0,+\infty)}^{}1_{(a,b]}d \nu_f = \nu_f(a,b] = \lambda_f(a) - \lambda_f(b) =  \int_{\Omega}^{}1_{(a,b]}(|f|) d \mu
\] 
this is true by math 245B
Note that
\[
    1_{(a,b]}(|f|) =
    \begin{cases}
        0 & |f| \not\in (a,b]\\
        1 & |f| \in (a,b]
    \end{cases}
    = 1_{|f|^{-1}((a,b])}
\] 
so
\[
    \int_{\Omega}^{}1_{(a,b]}(|f|)d \mu = \mu \{a < |f| \le b\}
\] 
Lemma: If $\lambda_f(\alpha) < +\infty$ for all $\alpha > 0$ then  for every $\phi: [0,+\infty) \to [0, +\infty)$ which is Borel measurable
\[
\int_{\Omega}^{}\phi(|f|) d \mu = \int_{(0,+\infty)}^{}\phi(\alpha) \nu_f(d \alpha) 
\] 
(Here we can replace $\phi \in L^{1}(\nu_f)$ by splitting into positive and negative parts)
Proof: It suffices to prove the lemma when $\phi = 1_{(a,b]}$ and $(a,b] \subset (0,+\infty)$ this is (1).
\[
\int_{\Omega}^{}\phi(|f|) d \mu = \int_{0}^{+\infty}\phi(\alpha) \nu_f(d \alpha) "=" \int_{0}^{+\infty} \phi(\alpha) (-\lambda_f(\alpha))' d \alpha =  \int_{0}^{\infty} \phi'(\alpha) \lambda_f(\alpha) d \alpha
\] 
which is saying that $\nu_f = -\lambda_f'(\alpha) d \alpha$ and we take $\phi = 0$ at 0 and infinity and use integration by parts.
If we can show this holds for all polynomial we can use an approximation argument to show the conclusion.
\section{Theorem}
For every $p > 0$, we have
\[
\int_{\Omega}^{}|f|^{p} d \mu = p \int_{0}^{+\infty}\alpha^{p-1}\lambda_f(\alpha) d \alpha
\] 
Proof: Case 1, there exists $\overline{\alpha} \in (0, +\infty)$ such that $\lambda_f(\overline{\alpha}) = +\infty$ Note that if $\alpha^{*} \in (0, \overline{ \alpha})$ then $\lambda_f(\alpha) = +\infty$ for $\alpha \in (\alpha^{*}, \overline{ \alpha})$
and so
\[
\int_{0}^{+\infty}\alpha^{p-1}\lambda_f(\alpha) d \alpha \ge \int_{\alpha^{*}}^{\overline{\alpha}} \alpha^{p-1} \lambda_f(\alpha) d \alpha = + \infty
\] 
and
\[
    \int_{\Omega}^{}|f|^{p}d \mu \ge \int_{\{\overline{\alpha} < |f|\}}^{}|f|^{p} d \mu \ge \overline{\alpha}^{p} \lambda_f( \overline{ \alpha) = +\infty}
\] 
Case 2: $\lambda_f(\alpha) < +\infty$ for all $\alpha \in (0, +\infty)$ let  $f = 1_A$ and then $f = \sum_{i=1}^{n}\lambda_i 1_{A_i}$ and then $f = \sum_{i=1}^{\infty}\lambda_i 1_{A_i}$ for the finite sum case assume that $A_i \cap A_j = \emptyset$ if  $i \ne j$.
 \[
     (2)    |f|^{p} = \sum_{}^{}\lambda_i^{p} 1_{A_i}
\] 
Observe that if $f$ is in this form $f = \sum_{i=1}^{n}f_i$ where $f_i$ are such that  $\{|f_i| > 0\} \cap \{|f_j| > 0\} = \emptyset$ for  $i \ne j$
then  $|f|^{p} = \sum_{i=1}^{n}|f_i|^{p}$ and $\lambda_f = \sum_{i=1}^{n} \lambda_{f_i}$.
Let us prove the result when $|f_i| = t_i1_{A_i}$ we have for $\alpha > 0$ and $t_i > 0$
 \[
     \lambda_{f_i}(\alpha) = \mu\{ \frac{\alpha}{t_i} < 1_{A_i}\} = 
     \begin{cases}
         0 & \alpha \ge t_i\\ \mu[A_i] & \alpha < t_i \end{cases}
\] 
Thus,
\[
    (3) \quad    p \int_{0}^{+\infty}\alpha^{p-1} \lambda_{f_i}(\alpha) d \alpha = p \int_{0}^{t_i} \alpha^{p-1} \mu[A_i] d \alpha = t_i^{p} \mu[A] = \int_{\Omega}^{}|f_i|^{p} d \mu
\] 
By $(2)$ and  $(3)$
 \[
\int_{\Omega}^{}|f|^{p} d\mu = \sum_{i=1}^{n} \int_{\Omega}^{}|f_i|^{p} d\mu = \sum_{i=1}^{n} p \int_{0}^{+\infty} \alpha^{p-1}\lambda_f(\alpha) d \alpha = p \int_{0}^{+\infty} \alpha^{p-1} \lambda_f(\alpha) d \alpha
\] 
Change of variables in $ \mathbb{R}^{d}$. Reminder: We call $O \in \mathbb{R}^{d \times d}$ orthogonal if $O^{T}O = I_d = OO^{T}$ 
If $S \in \mathbb{R}^{d \times d}$ is symmetric then there exists $O \in \mathbb{R}^{d\times d}$ orthogonal and $\Lambda = diag(\lambda_1,\dots,\lambda_d)$ such that
\[
S = O^{T} \Lambda O
\] 
If $A \in \mathbb{R}^{d \times d}$ then there exists $P \in \mathbb{R}^{d \times d}$ orthogonal and $S \in \mathbb{R}^{d \times d}$ symmetric such that $A = PS$.
We identiyfy $A$ with  $x \to Ax$
Set 
\[
    B_r(a) = \{x \in \mathbb{R}^{d} : ||x -a||_{l_2} \le r\} \qquad Q_r(a) = [a_1-r,a_1+r] \times \dots \times [a_d - r, a_d + r]
\] 
Note that
\[
O(B_r(a)) = B_r(O(A))
\] 
\[
    \Lambda Q_r(a) = [\lambda_1 a_1 - r |\lambda_1|, \lambda_1 a_1 + r|\lambda_1|] \times \dots \times [\lambda_d a_d - r |\lambda_d|, \lambda_d a_d + r |\lambda_d|]
\] 
\[
L^{d}(B_r(a)) = L^{d}(O(B_r(a)) \implies L^{d}(E) = L^{d}(O(E))
\] 
and
\[
L^{d}(\Lambda Q_r(a)) = (2r |\lambda_1|) \dots (2r |\lambda_d|) = |det \Lambda| L^{d}(Q_r(a)) 
\] 
\[
\implies L^{d}(\Lambda E) = |det \Lambda| L^{d}(E)
\] 

\section{Friday}
Fact(1) From Homework 2 we derive the fact: Let $A \subset \mathbb{R}^{d}$ be $\mathcal{L}^{d}$ measurable.
Then for every $\epsilon > 0$ there exist balls $(B_{r_i}(x_i))_{i=1}^{\infty}$ such that
\[
    A \subset \bigcup_{i=1}^{\infty}B_{r_i}(x_i) \quad \text{ and } \mathcal{ L}^{d}(A) \ge \sum_{i=1}^{\infty} \mathcal{L}^{d}(B_{r_i}(x_i)) - \epsilon
\] 
Given $\delta > 0$, $O \subset \mathbb{R}^{d}$ open, there exists $(B_{r_i}(x_i))_{i=1}^{\infty}$ disjoint such that
\[
    r_i \le \delta \quad B_{r_i}(x_i) \subset O \qquad \mathcal{L}(A \cap O \setminus \bigcup_{i=1}^{\infty} B_{r_i}(x_i)) = 0
\] 

We are going to assume the following facts:
Fact(2)
\[
\mathcal{L}^{d}( B_r(x_0)) = r^{d} \mathcal{L}^{d}(B_1( \vec 0))
\] 
Exercise: (Assume we can solve EXO 2.1 and 2.2)\\
Let $M \in \mathbb{R}^{d \times d}$ be a matrix and $A \subset \mathbb{R}^{d}$ be Borel. Show that $ \mathcal{L}^{d}( \{ Mx : x \in A\}) = | det M| \mathcal{L}^{d}(A)$.
Identify $M$ with  $x \to Mx$ and assume we take
\[
f : \mathbb{R}^{d} \to \mathbb{R}^{d} \quad \mathcal{L}^{d}(f(A))
\] 
Solution: By Facts 1, it suffices to solve the exercise when $A = B_r(x_0)$ write $M = PS$ where  $P, S \in \mathbb{R}^{d \times d}$ $P$ orthogonal,  $S$ symmetric
and write  $S = Q^{\top} \Lambda Q$ where $Q, \Lambda \in \mathbb{R}^{d}$ and $Q$ orthogonal,  $\Lambda = diag(\lambda_1,\dots,\lambda_d)$ then
\[
\mathcal{L}^{d}(M B_r(x_0)) = \mathcal{L}^{d}( P Q^{\top} \Lambda Q B_r(x_0)) = \mathcal{L}^{d}(\Lambda Q B_r(x_0) = \mathcal{L}^{d}( \Lambda B_r(Q(x_0))) = |det \Lambda| \mathcal{L}^{d}(B_r(Q(x_0))) 
\] 
but $|det \Lambda| = |det M|$.
Goal: Let  $f : \mathbb{R}^{d} \to \mathbb{R}^{d}$ be Lipschitz and let $g \in  L^{1}(\mathbb{R}^{d})$. Set $J_f = det \nabla f$ where we use Rademacher theorem that says the derivative
exists almost everywhere. We want to show that
 \[
\int_{ \mathbb{R}^{d}}^{}g(x) |J_f(x)| dx  = \int_{ \mathbb{R}^{d}}^{} \sum_{x \in f^{-1}(y)}^{}g(x) dy
\] 
In particular if $\tilde g : \mathbb{R}^{d} \to [0,+\infty)$ is Borel then $g = \tilde g \circ f$ and  $D \subset \mathbb{R}^{d}$ is open then
\[
\int_{ D}^{} \tilde g \circ f |J_f(x)| dx = \int_{ \mathbb{R}^{d}}^{}\tilde g(y) N(y,D) dy
\] 
where $N(y,D) = Card\{f^{-1}(y) \cap D\}$ and $y = \nabla f(x), dy = |det \nabla f(x)|dx$\\
Exercise: Let  $f : \mathbb{R}^{d} \to \mathbb{R}^{d}$ be $l$-Lipshchitz and let  $A \subset \mathbb{R}^{d}$ be  $ \mathcal{L}^{d}$-measurable.
\begin{enumerate}
    \item Show that $f(A) $ is  $ \mathcal{L}^{d}$-measurable
    \item Show that $y \to N(y,A) = Card\{f^{-1}(y) \cap A\}$ is $ \mathcal{L}^{d}$ - measurable.
\end{enumerate}
Solution: It is a not a loss of generality to assume that $A$ is bounded.
\begin{enumerate}
    \item For each $n \in \mathbb{N}$ we choose $K_n \subset A$ compact such that $ \mathcal{L}^{d}(A \setminus K_n) < \frac{1}{n}$. Since $f(K_n)$ is compact,  $\bigcup_{n=1}^{\infty}f(K_n)$ is a borel set. Note that $ \mathcal{L}^{d}(A \setminus ( \bigcup_{m=1}^{\infty}K_m)) \le \mathcal{L}^{d}(A \setminus K_n) \le \frac{1}{n}$ so
        it is a null set $ \mathcal{L}^{d}(A \setminus (\bigcup_{m=1}^{\infty}K_m)) = 0$. Let $K = bigcup_{m=1}^{\infty}K_m$ 
        \[
            f(A) = f(K) \cup f(A \setminus K)
        \] 
        \[
        \mathcal{L}^{d}(f(A \setminus K)) \le (Lip f)^{d} \mathcal{L}^{d}(A \setminus K)  = 0
        \] 
        by exercise $2.3$ 
        Hence $f(A \setminus K)$ is  $ \mathcal{L}^{d}$ measurable and so $f(A)$ is  $ \mathcal{L}^{d}$-measurable.
    \item 
        I am going to divide $ \mathbb{R}^{d}$ into rectangles dyadic and in each rectangle given $y$ we want to count the number of  $x$ such that  $f(x) = y$ if there are finitely many points in the preimage we can make the grdid so small that each are in a separate grid, otherwise there are infinite and( i lost what he said).
        Fore each  $k \in \mathbb{N}$ set  $B_k = \{Q : Q = (a_1,b_1] \times \dots \times (a_d,b_d] : b_i = a_i + \frac{1}{k}, a_i \in \frac{1}{k} \mathbb{Z}, i = 1,\dots,d\}$.
        Note that $ B_k$ is a partition of $ \mathbb{R}^{d}$. 
        givne $y \in \mathbb{R}^{d}$ we set
        \[
            g_k(y) = Card\{Q : Q \in B_k : f^{-1}(y) \cap (A \cap Q) \ne \emptyset\} = \sum_{Q \in B_k}^{}1_{f(A \cap Q)}
        \] 
        By i, $g_k$ is  $ \mathcal{L}^{d}$ measurable and by homework 2.4 or 2.5
        \[
            \lim_{k\to \infty}g_k(y) = card\{ f^{-1}(y) \cap A\} = N(y,A)
        \] 
        Remark, By the monotone convergence theorem 
        \[
        \int_{ \mathbb{R}^{d}}^{}N(y,A) dy = \lim_{k\to \infty} \int_{ \mathbb{R}^{d}}^{}g_k(y)dy = \lim_{k\to \infty} \sum_{ Q \in B_k}^{} \mathcal{L}^{d}(f^{-1}(A \cap Q))\]\[ \le \lim_{k\to \infty} \sum_{ Q \in B_k}^{} (Lip f)^{d} \mathcal{L}^{d}(A \cap Q) = (Lip f)^{d} \mathcal{L}^{d}(A)
        \] 
        Hence $N(\cdot, A) \in L^{1}( \mathbb{R}^{d})$ if $A$ is bounded
\end{enumerate}




\end{document}
